% ============================================================
%  DADOS DO TRABALHO
%  Preencha APENAS este arquivo com suas informacoes.
%  Nao mexa nos outros arquivos.
% ============================================================

% --- SEU NOME ---
\newcommand{\oautor}{Seu Nome Completo Aqui}
% Exemplo: \newcommand{\oautor}{Maria Clara Souza Oliveira}

% --- NOME DO ORIENTADOR (com titulo) ---
\newcommand{\oorientador}{Prof. Dr. Nome do Orientador}
% Exemplos: Prof. Dr. / Prof. Me. / Prof. Esp.

% --- TITULO DO TRABALHO ---
\newcommand{\otitulo}{Titulo do Seu Projeto de Pesquisa}
% Exemplo: \newcommand{\otitulo}{Sistema Web para Gestao de
%          Filas em Postos de Saude Municipais}

% --- SUBTITULO (deixe vazio se nao tiver) ---
\newcommand{\osubtitulo}{}
% Com subtitulo: \newcommand{\osubtitulo}{Uma Abordagem com Microsservicos}

% --- TIPO DO TRABALHO ---
\newcommand{\otipotrabalho}{Projeto de Pesquisa}
% PS I  -> Projeto de Pesquisa
% PS II -> Trabalho de Conclusao de Curso

% --- LOCAL E ANO ---
\newcommand{\olocal}{Tiangua -- CE}
\newcommand{\oano}{2026}

% --- PALAVRAS-CHAVE (3 a 6, separadas por ponto) ---
\newcommand{\opalavraschave}{%
  Palavra-chave 1. Palavra-chave 2. Palavra-chave 3. Sistemas de Informacao.}
% Exemplo:
% \newcommand{\opalavraschave}{%
%   Sistema Web. Gestao de Filas. Postos de Saude. Sistemas de Informacao.}

\newcommand{\okeywords}{%
  Keyword 1. Keyword 2. Keyword 3. Information Systems.}

% --- TEXTO DO PREAMBULO (nao altere) ---
\newcommand{\opreambulo}{%
  \otipotrabalho{} apresentado a disciplina de
  Projeto Supervisionado ou de Graduacao I do curso
  de Bacharelado em Sistemas de Informacao da
  Faculdade UNINTA, como requisito parcial para
  a elaboracao do Trabalho de Conclusao de Curso.}
